\documentclass[11pt]{article}

\usepackage[margin=1in]{geometry}
\usepackage{ctex}
\usepackage{booktabs}
\usepackage{array}

\title{InfoGate 框架动机简述（中文版）}
\author{}
\date{2026年4月11日}

\begin{document}

\maketitle

\section{背景}

多模态情感分析需融合语言、声学、视觉信号。难点不仅是``如何融合''，而是\textbf{何者主导当前样本的决策}，以及如何避免次要模态中的噪声污染主信息流。对称融合会稀释最强信号；固定以语言为中心在声学与视觉起决定作用的样本上会失效。因此\textbf{按样本动态选择主模态}是合理方向，但若表征仍含任务无关变异、序列冗余或不可靠跨模态证据，则路由与注入都会不稳定。

InfoGate 针对这一缺口：先在更任务相关的空间中\textbf{净化}多模态表征，再在该空间内做动态主模态选择，最后以\textbf{可靠性感知}的方式向主流注入辅助信息。核心可概括为：\textbf{噪声感知、样本自适应、以主模态为中心的融合}。

\section{相关工作（提要）}

\begin{itemize}
  \item \textbf{对称融合}：三模态近似对等交互，隐含``各模态同等可信''，易混入干扰。
  \item \textbf{静态语言中心}：全局先验合理，但在情感主要由声调或表情表达的样本上失效。
  \item \textbf{动态主模态（如 MODS）}：主模态随样本而变，但若仍用原始或弱过滤特征，噪声问题未根本解决。
  \item \textbf{信息瓶颈与信息空间（ITHP、CyIN）}：将多模态压入更含任务信息、抑制干扰的潜空间；InfoGate 将此类思想与主模态中心融合结合。
\end{itemize}

\section{三条动机}

\begin{enumerate}
  \item \textbf{路由需要干净表征}：若声学/视觉仍冗余、文本仍含与情感无关的词面变化，主模态选择会不稳定；选择器应基于\textbf{面向任务的特征}而非原始投影。
  \item \textbf{辅助模态并非同等有用}：标准交叉注意力假设辅助 token 均可平等影响主流，过强；需反映不确定性，\textbf{衰减}不可靠的跨模态证据。
  \item \textbf{预测应保持主模态中心}：若末端仍大量拼接各模态残差，分类器可能绕开既定路由；更合理的是仅从\textbf{增强后的主瓶颈流}输出预测。
\end{enumerate}

\section{问题与对策（简表）}

\begin{table}[h]
\centering
\begin{tabular}{>{\raggedright\arraybackslash}p{0.28\linewidth} >{\raggedright\arraybackslash}p{0.62\linewidth}}
\toprule
问题 & 当前框架中的对策 \\
\midrule
样本间主模态不同 & MSelector 按样本动态选主模态，而非固定以语言为中心。 \\
单模态特征嘈杂、含无关信息 & 各模态经信息瓶颈编码，路由与融合在压缩潜表征上进行。 \\
辅助模态可靠性不均 & 由瓶颈不确定性得到 token 级置信度，调制交叉注意力。 \\
辅助信息不应同等注入 & 自适应信息门控制各辅助分支对主流的贡献。 \\
分类可能绕开主模态融合 & 取消直接多模态捷径，仅从增强主瓶颈表征预测。 \\
跨模态语义在训练中漂移 & 瓶颈重构与类 CRA 的翻译正则，使交互留在一致的信息空间内。 \\
\bottomrule
\end{tabular}
\caption{InfoGate 的问题--对策对应关系（简述）。}
\end{table}

\section{流程一句话}

各模态经共享隐空间后进入信息瓶颈，得到压缩特征与置信度；池化后由 MSelector 路由主模态；两路辅助经置信度调制的交叉注意力与门控注入主流；最终仅聚合增强主模态用于预测。

\section{讨论与一句总结}

叙事主线清晰：\textbf{先过滤信息，再选主模态，再在不确定性控制下融合辅助，最后从增强主表征决策}。表述上宜强调完整三模态场景与动态主模态、噪声感知融合；若未重新评测，不宜过度宣称缺失模态设定。翻译式正则若保留于代码，亦不应单独支撑``缺失模态系统''的强结论。

\textbf{一句话：}已有主模态中心方法承认不同样本需要不同主导模态，但对进入选择与融合的信息质量与可靠性控制不足；本框架以信息瓶颈引导、置信度感知、动态路由的主模态中心结构，使模态选择与跨模态交互更稳定、更贴近任务。

\end{document}
